\documentclass[11pt]{article}
\usepackage[T1]{fontenc}
\usepackage[spanish]{babel}
\usepackage{amsmath,amssymb}
\usepackage[a4paper,margin=2.3cm]{geometry}
\usepackage{graphicx}
\usepackage{tikz}
\usetikzlibrary{arrows.meta,patterns.meta}

\IfFileExists{densidad_plano_visualizacion.pdf}{%
  \newcommand{\includevisualizacion}[2][]{\includegraphics[##1]{##2.pdf}}%
}{%
  \usepackage[inkscapelatex=false]{svg}
  \newcommand{\includevisualizacion}[2][]{\includesvg[##1]{##2}}%
}

\spanishdecimal{.}
\setlength{\parindent}{0pt}
\setlength{\parskip}{0.35em}
\setlength{\abovedisplayskip}{4pt}
\setlength{\belowdisplayskip}{4pt}
\setlength{\abovedisplayshortskip}{4pt}
\setlength{\belowdisplayshortskip}{4pt}

\newcommand{\encabezado}[3]{%
  {\Large\bfseries #1\par}
  \vspace{2pt}
  {\normalsize #2\hfill #3\par}
  \vspace{6pt}\hrule\vspace{8pt}}

\newcommand{\problema}[2]{%
  \vspace{0.55em}
  \noindent\textbf{#1.}\enspace{\small #2}\par
  \vspace{0.3em}}

\begin{document}

\encabezado{Densidad superficial de carga a partir del potencial}
{Electrodinámica I}{P.S., J.A., J.R.}

\problema{1}{El plano infinito $z=0$ divide el vacío en dos regiones. El
potencial eléctrico está dado por
\begin{equation}
\phi(x,y,z)=
\begin{cases}
\displaystyle \frac{V_0z}{(x^2+y^2+z^2)^{3/2}}, & z>0,\\[8pt]
\displaystyle \frac{V_0z}{2(x^2+y^2+z^2)^{3/2}}, & z<0,
\end{cases}
\label{eq:potencial-plano}
\end{equation}
donde $V_0>0$ y $[V_0]=[\phi]L^2$. Determine:
\begin{enumerate}
  \item el campo eléctrico en las regiones superior e inferior;
  \item sus componentes tangencial y normal al aproximarse a $z=0$;
  \item la continuidad de la componente tangencial;
  \item la densidad superficial de carga $\sigma(x,y)$ sobre el plano.
\end{enumerate}
Todos los límites sobre la interfaz se entienden para
$\rho=\sqrt{x^2+y^2}>0$, pues el potencial dado es singular en el origen.}

\begin{figure}[htbp]
\centering
\begin{tikzpicture}[
    scale=0.93,
    line cap=round,
    line join=round,
    eje/.style={-{Stealth[length=1.9mm]},semithick},
    vector/.style={-{Stealth[length=1.9mm]},thick},
    every node/.style={font=\small}
]
\begin{scope}[x={(1cm,0cm)},y={(-0.44cm,0.24cm)},z={(0cm,1cm)}]
  \path[
    pattern={Lines[angle=35,distance=4.2pt,line width=0.22pt]},
    pattern color=black!42
  ] (-3.25,-2.20,0)--(3.25,-2.20,0)--(3.25,2.20,0)--(-3.25,2.20,0)--cycle;
  \draw[thick] (-3.25,-2.20,0)--(3.25,-2.20,0)--(3.25,2.20,0)--(-3.25,2.20,0)--cycle;

  \coordinate (O) at (0,0,0);
  % La proyección del eje y apunta hacia la izquierda; así la normal y el
  % punto P disponen de una zona propia en el costado derecho del dibujo.
  \coordinate (S) at (2.00,0.70,0);
  \draw[eje] (-3.75,0,0)--(3.85,0,0) node[below] {$x$};
  \draw[eje] (0,-3.10,0)--(0,3.15,0) node[above right=-1pt] {$y$};
  \draw[eje] (0,0,-1.55)--(0,0,2.95) node[left] {$z$};
  \fill (O) circle (1.1pt);
  \node[below left] at (O) {$O$};
  \fill (S) circle (1.6pt);
  \node[anchor=west,fill=white,inner sep=1.2pt]
    at (2.14,0.70,0) {$P$};
  \draw[vector] (O)--(S)
    node[pos=.56,below=4pt,fill=white,inner sep=1pt] {$\vec\rho$};
  \draw[vector] (S)--++(0,0,1.45)
    node[midway,right=4pt,fill=white,inner sep=1pt] {$\hat n=\hat z$};

  \node[anchor=west,align=left,fill=white,inner sep=2pt]
    at (3.05,0,2.02)
    {$z>0$\quad $\phi^+=V_0z/r^3$};
  \node[anchor=west,align=left,fill=white,inner sep=2pt]
    at (3.05,0,-1.22)
    {$z<0$\quad $\phi^-=V_0z/(2r^3)$};
  \node[anchor=north,fill=white,inner sep=2pt]
    at (-1.70,1.60,0) {$z=0$};
\end{scope}
\end{tikzpicture}

\caption{El punto $P=(x,y,0)$ se localiza mediante
$\vec\rho=x\hat x+y\hat y$, con $\rho=\sqrt{x^2+y^2}$. La normal
$\hat n=\hat z$ apunta hacia la región superior y
$r=\sqrt{x^2+y^2+z^2}$.}
\label{fig:geometria-plano}
\end{figure}

\textbf{Solución:}

\textbf{1. Cálculo del campo eléctrico por regiones.}\par
El campo se obtiene de la relación
\begin{equation}
\vec E=-\nabla\phi
=-\left(
\frac{\partial\phi}{\partial x}\,\hat x
+\frac{\partial\phi}{\partial y}\,\hat y
+\frac{\partial\phi}{\partial z}\,\hat z
\right).
\label{eq:campo-gradiente}
\end{equation}
Escribimos
\[
r=\sqrt{x^2+y^2+z^2}.
\]
En la región superior, $\phi^+=V_0zr^{-3}$. Sus derivadas son
\begin{equation*}
\frac{\partial\phi^+}{\partial x}=-\frac{3V_0xz}{r^5},
\qquad
\frac{\partial\phi^+}{\partial y}=-\frac{3V_0yz}{r^5},
\qquad
\frac{\partial\phi^+}{\partial z}
=\frac{V_0(r^2-3z^2)}{r^5}.
\end{equation*}
Por lo tanto,
\begin{equation}
\boxed{
\vec E^{+}(x,y,z)=\frac{V_0}{r^5}
\left[3xz\,\hat x+3yz\,\hat y+(3z^2-r^2)\,\hat z\right],
\qquad z>0.}
\label{eq:campo-superior}
\end{equation}

En la región inferior, el potencial tiene exactamente la misma dependencia
espacial, multiplicada por $1/2$. En consecuencia,
\begin{equation}
\boxed{
\vec E^{-}(x,y,z)=\frac{V_0}{2r^5}
\left[3xz\,\hat x+3yz\,\hat y+(3z^2-r^2)\,\hat z\right],
\qquad z<0.}
\label{eq:campo-inferior}
\end{equation}

\textbf{2. Evaluación sobre la frontera.}\par
Para aproximarnos al plano fijamos $(x,y)$ y tomamos $z\to0^\pm$. Como
$r\to\rho=\sqrt{x^2+y^2}$, los términos proporcionales a $xz$, $yz$ y $z^2$
se anulan. De las ecuaciones~\eqref{eq:campo-superior} y
\eqref{eq:campo-inferior} obtenemos
\begin{equation}
\boxed{
\vec E^{+}(x,y,0)=-\frac{V_0}{\rho^3}\,\hat z,
\qquad
\vec E^{-}(x,y,0)=-\frac{V_0}{2\rho^3}\,\hat z,}
\qquad \rho>0.
\label{eq:campos-frontera}
\end{equation}

Respecto de la normal $\hat n=\hat z$, definimos
\[
E_\perp=\vec E\cdot\hat n,
\qquad
\vec E_\parallel=\vec E-E_\perp\hat n.
\]
Así,
\begin{equation}
E^+_\perp=-\frac{V_0}{\rho^3},
\qquad
E^-_\perp=-\frac{V_0}{2\rho^3},
\qquad
\vec E^+_\parallel=\vec E^-_\parallel=\vec 0.
\label{eq:componentes-frontera}
\end{equation}

\begin{figure}[htbp]
\centering
\begin{tikzpicture}[
    scale=0.88,
    line cap=round,
    line join=round,
    vector/.style={-{Stealth[length=1.9mm]},thick},
    recorrido/.style={-{Stealth[length=1.7mm]},semithick},
    every node/.style={font=\small}
]
% Panel (a): contorno de circulación.
\begin{scope}[xshift=-4.05cm]
  \node[font=\small\bfseries] at (0,1.78) {(a) Contorno de circulación};
  \draw[very thick] (-3.05,0)--(3.05,0);
  \node[above right] at (2.25,0) {$z=0$};
  \node[anchor=west] at (-2.90,1.33) {región $+$};
  \node[anchor=west] at (-2.90,-1.33) {región $-$};

  \draw[recorrido] (-1.50,0.62)--(1.50,0.62);
  \draw[recorrido] (1.50,0.62)--(1.50,-0.62);
  \draw[recorrido] (1.50,-0.62)--(-1.50,-0.62);
  \draw[recorrido] (-1.50,-0.62)--(-1.50,0.62);
  \node[above] at (0,0.62) {$\Delta\ell$};
  \draw[{Stealth[length=1.4mm]}-{Stealth[length=1.4mm]}]
    (1.94,-0.60)--(1.94,0.60);
  \node[right] at (1.94,0) {$2\delta$};
  \draw[vector] (-0.88,1.17)--(0.30,1.17)
    node[right] {$\vec E_\parallel^+$};
  \draw[vector] (0.30,-1.17)--(-0.88,-1.17)
    node[left] {$\vec E_\parallel^-$};
  \draw[vector] (-2.35,0.20)--(-1.35,0.20)
    node[above] {$\hat t$};
\end{scope}

% Panel (b): pastilla gaussiana.
\begin{scope}[xshift=4.05cm]
  \node[font=\small\bfseries] at (0,1.78) {(b) Pastilla gaussiana};
  \draw[very thick] (-3.05,0)--(3.05,0);
  \node[above right] at (2.25,0) {$z=0$};
  \node[anchor=west] at (-2.90,1.33) {región $+$};
  \node[anchor=west] at (-2.90,-1.33) {región $-$};

  \draw[densely dashed] (-1.15,0.62) ellipse (1.15 and 0.23);
  \draw[densely dashed] (-1.15,-0.62) ellipse (1.15 and 0.23);
  \draw[densely dashed] (-2.30,0.62)--(-2.30,-0.62);
  \draw[densely dashed] (0,0.62)--(0,-0.62);
  \draw[vector] (-1.40,1.36)--(-1.40,0.80)
    node[midway,left=2pt] {$\vec E^+$};
  \draw[vector] (-1.40,-0.80)--(-1.40,-1.32)
    node[midway,left=2pt] {$\vec E^-$};
  \draw[vector] (0,0.66)--(0,1.34)
    node[right] {$+\hat n$};
  \draw[vector] (0,-0.66)--(0,-1.34)
    node[right] {$-\hat n$};
  \foreach \x in {-2.65,-2.10,-1.55,-1.00,-0.45,0.10,0.65,1.20,1.75,2.30,2.85}{
    \node[font=\scriptsize] at (\x,0) {$-$};
  }
  \node[below=3pt,fill=white,inner sep=1pt] at (1.36,0) {$\sigma(\rho)<0$};
\end{scope}
\end{tikzpicture}
\caption{Geometrías usadas en las condiciones de frontera. El contorno de (a)
compara las componentes tangenciales; la pastilla de (b) relaciona el salto de
la componente normal con la densidad superficial.}
\label{fig:frontera-plano}
\end{figure}

\textbf{3. Continuidad de la componente tangencial.}\par
Dividimos el contorno de la Figura~\ref{fig:frontera-plano}(a) en sus cuatro
lados,
\[
C=C_1\cup C_2\cup C_3\cup C_4.
\]
Por el teorema de Stokes y $\nabla\times\vec E=\vec0$,
\begin{align*}
0=\oint_C\vec E\cdot d\vec\ell
&=\int_{C_1}\vec E\cdot d\vec\ell
+\int_{C_2}\vec E\cdot d\vec\ell\\
&\quad+\int_{C_3}\vec E\cdot d\vec\ell
+\int_{C_4}\vec E\cdot d\vec\ell.
\end{align*}
Al hacer $\delta\to0$, los lados cortos desaparecen porque su longitud tiende a
cero. Los lados largos entregan
\[
\int_{C_1}\vec E\cdot d\vec\ell
\longrightarrow \vec E_\parallel^+\cdot\hat t\,\Delta\ell,
\qquad
\int_{C_3}\vec E\cdot d\vec\ell
\longrightarrow-\vec E_\parallel^-\cdot\hat t\,\Delta\ell.
\]
El signo menos aparece porque el lado inferior se recorre en el sentido
contrario. Por tanto,
\[
0=\left(\vec E_\parallel^+-\vec E_\parallel^-\right)
\cdot\hat t\,\Delta\ell,
\]
y, para cualquier dirección tangente $\hat t$,
\begin{equation}
\vec E^+_\parallel-\vec E^-_\parallel=\vec 0.
\label{eq:continuidad-tangencial}
\end{equation}
En este problema la verificación es inmediata a partir de
\eqref{eq:componentes-frontera}: ambas componentes tangenciales son nulas. El
potencial también es continuo, pues
\[
\lim_{z\to0^+}\phi(x,y,z)
=\lim_{z\to0^-}\phi(x,y,z)=0,
\qquad \rho>0.
\]

\textbf{4. Densidad superficial de carga.}\par
Dividimos la superficie de la Figura~\ref{fig:frontera-plano}(b) en tapa
superior, pared lateral y tapa inferior:
\[
S=S_+\cup S_{\mathrm{lat}}\cup S_-.
\]
La Ley de Gauss entrega
\begin{align*}
\frac{q_{\mathrm{enc}}}{\varepsilon_0}
&=\oint_S\vec E\cdot d\vec S\\
&=\int_{S_+}\vec E\cdot d\vec S
+\int_{S_{\mathrm{lat}}}\vec E\cdot d\vec S
+\int_{S_-}\vec E\cdot d\vec S.
\end{align*}
Al hacer $\delta\to0$,
\[
\int_{S_+}\vec E\cdot d\vec S\longrightarrow E_\perp^+A,
\qquad
\int_{S_-}\vec E\cdot d\vec S\longrightarrow-E_\perp^-A,
\qquad
\int_{S_{\mathrm{lat}}}\vec E\cdot d\vec S\longrightarrow0.
\]
El flujo lateral desaparece porque el área de la pared tiende a cero. Como la
pastilla encierra $q_{\mathrm{enc}}=\sigma A$,
\[
\left(E_\perp^+-E_\perp^-\right)A
=\frac{\sigma A}{\varepsilon_0}.
\]
Cancelando el área, obtenemos
\begin{equation}
\hat n\cdot\left(\vec E^+-\vec E^-\right)
=\frac{\sigma}{\varepsilon_0}.
\label{eq:salto-normal}
\end{equation}
El signo menos de $E_\perp^-$ se debe a que la normal exterior de la tapa
inferior es $-\hat n$.

Sustituyendo la ecuación~\eqref{eq:campos-frontera},
\begin{align*}
\frac{\sigma}{\varepsilon_0}
&=-\frac{V_0}{\rho^3}-\left(-\frac{V_0}{2\rho^3}\right)\\
&=-\frac{V_0}{2\rho^3}.
\end{align*}
Por tanto,
\begin{equation}
\boxed{
\sigma(x,y)=-\frac{\varepsilon_0V_0}
{2(x^2+y^2)^{3/2}},}
\qquad (x,y)\ne(0,0).
\label{eq:sigma-final}
\end{equation}

La densidad es negativa en todo punto regular del plano y su magnitud decrece
como $\rho^{-3}$. La divergencia cuando $\rho\to0$ no es un defecto del cálculo:
refleja la singularidad del potencial impuesto en el origen. Un modelo físico
con tamaño finito debe regularizar esa zona; por ejemplo, al excluir un disco de
radio $a$ la carga contenida en $\rho>a$ es
\[
Q(\rho>a)=\int_a^\infty\sigma(\rho)\,2\pi\rho\,d\rho
=-\frac{\pi\varepsilon_0V_0}{a},
\]
que diverge en el límite ideal $a\to0$.

\section*{Visualización numérica}

Tomamos una longitud de referencia $L$ y representamos
\[
\widetilde{\sigma}=\frac{\sigma L^3}{\varepsilon_0V_0},
\qquad
\widetilde{\vec E}=\frac{\vec E L^3}{V_0}.
\]
La región central $\rho/L<0.6$ se deja sin representar porque el potencial
y la densidad son singulares en el origen. Se muestran tanto el perfil radial
y el mapa de $\widetilde{\sigma}$ como la dirección de
$\widetilde{\vec E}$ en un corte meridiano. La comparación de los límites
normales a ambos lados del plano permite comprobar directamente el salto.

\begin{figure}[htbp]
\centering
\includevisualizacion[width=0.97\textwidth]{densidad_plano_visualizacion}
\caption{Visualización numérica de las condiciones de frontera. (a) Perfil
radial de la densidad superficial. (b) Mapa de $\widetilde{\sigma}$ en
$z=0$, con el origen singular excluido. (c) Dirección del campo eléctrico en
el corte $y=0$; los colores distinguen las dos regiones. (d) Límites
$\widetilde E_z^+$ y $\widetilde E_z^-$ sobre el plano: su diferencia es
$\widetilde{\sigma}$.}
\label{fig:visualizacion-plano}
\end{figure}

\end{document}
